
\chapter{UnionPay}\label{ch:unionpay}
В~международном приёме UnionPay объединяет инфраструктуру UnionPay
International (UPI, не~путать с~индийской Unified Payments Interface)~\cite{unionpay-international-about}, правила конкретной
юрисдикции~\cite{unionpay-russia-rules} и~кобейджинговые маршруты; путь
по~кобейджу выбирается по~правилам соответствующей платёжной
системы.

\section{Процессинговая архитектура UnionPay}\label{sec:up-cups}

Внутренняя расчётная и~процессинговая инфраструктура China UnionPay (CUP)
обрабатывает авторизацию, клиринг и~расчёт карточных транзакций внутри Китая
и~выполняет в~UnionPay роль, сопоставимую с~VisaNet у~Visa. UPI не публикует
подробного открытого описания этой инфраструктуры; в~отраслевых материалах
используется обобщённая аббревиатура CUP, иногда с~уточнением назначения
подсистем: процессинг, клиринг, расчёт~\cite{unionpay-russia-rules}.

\highlight{Авторизация и~клиринг разнесены во~времени и~по~требованиям
к~задержкам так~же, как у~Visa и~Mastercard} (см.~гл.~\ref{ch:tx-lifecycle});
\highlight{внутренняя инфраструктура UnionPay маршрутизирует авторизационные
  запросы к~эмитенту, а~клиринговая обработка собирает исходящие реестры
  за~период, рассчитывает нетто-позиции и~передаёт их в~расчётный
центр}~\cite{unionpay-russia-rules}.

Опубликованные правила UnionPay в~России задают регламент обработки:
авторизации работают круглосуточно; исходящие реестры и~отчётность поступают
до~09:00 МСК дня~T; реестр нетто-позиций направляется в~расчётный центр
до~12:00 МСК дня~T+1; расчёты между участниками завершаются до~13:00 МСК
дня~T+1~\cite{unionpay-russia-rules}.\footnote{Внутрироссийские операции
  по~картам международных платёжных систем по~ч.~4 ст.~30.6 Федерального закона
  от~27.06.2011~№~161-ФЗ обрабатываются через операционный платёжный
  и~клиринговый центр НСПК~\cite{fz-161}; значения
сверяются по~актуальной редакции правил.}

Правила UnionPay в~России фиксируют сроки, но не определяют исполнителя.
Внутрироссийские операции по~картам UnionPay по~ч.~4 ст.~30.6 Федерального
закона от~27.06.2011~№~161-ФЗ
обрабатывает операционный платёжный и~клиринговый центр НСПК (ОПКЦ,
см.~раздел~\ref{sec:mir-opkc})~\cite{fz-161}; UPI остаётся правообладателем
товарного знака и~получает нетто-позиции по~трансграничной части через свой
расчётный центр, а~внутренний клиринг между российскими участниками проходит через
ОПКЦ. Регламент 09:00/12:00/13:00~МСК из~правил UnionPay применяется с~поправкой:
к~отчётным срокам ОПКЦ применяются графики НСПК, согласованные с~UnionPay,
и~участник сверяет оба календаря одновременно. При~расхождении между расчётным
днём UPI и~операционным днём НСПК для~российских участников определяющим
становится календарь НСПК; для~внешнего участника~--- календарь UPI.

\importantnote{%
  За~пределами материкового Китая операции проходят через инфраструктуру
  UnionPay International~--- связующее звено между участником в~стране приёма
  и~внутренней инфраструктурой UnionPay в~Китае. Каждая
  юрисдикция применяет собственную редакцию правил и~регламента;
  сроки обработки и~форматы файлов проверяются отдельно
  по~каждой стране.
}

\section{Операционный цикл и~расчётный регламент}\label{sec:up-architecture}

Открытые правила UnionPay в~России описывают двухэтапную карточную схему.
Эквайер направляет эмитенту авторизационный запрос через операционный центр;
после одобрения эквайер отправляет исходящий клиринговый файл; платёжный
клиринговый центр рассчитывает нетто-позиции; расчётный центр исполняет их
по~счетам участников~\cite{unionpay-russia-rules}. По~умолчанию UnionPay
обрабатывает транзакцию двумя сообщениями (\textit{dual-message}), а~не одним
(\textit{single-message}).

В~отличие от~Visa BASE~II и~Mastercard GCMS, расчётный цикл UnionPay
изначально строился под~внутрикитайский рынок, а~зарубежный приём
обслуживается через UnionPay International. Внутри Китая инфраструктура
UnionPay обрабатывает транзакции с~расчётами T+1 по~умолчанию;
для~международных операций схема зависит от~соглашения между эквайером
и~эмитентом через UPI. Сроки расчёта по~одной и~той~же карте UnionPay в~разных
юрисдикциях сверяются по~актуальной редакции правил соответствующей расчётной
системы~\cite{unionpay-russia-rules}.

Single-message операции занимают в~UnionPay более важное место, чем в~Visa
или Mastercard. UnionPay вырос из~дебетового и~банкоматного рынка Китая
\mbox{2002--2010}~годов, где транзакция сводилась к~снятию наличных в~ATM и~покупке
с~моментальным списанием. Для~дебетовой карты, где нет кредитного лимита и~нет
разрыва между авторизацией и~расчётом, single-message естественнее: одно
сообщение содержит и~авторизацию, и~финансовый факт. Visa и~Mastercard, напротив,
развивались из~кредитной модели, где холд, подтверждение и~расчёт разнесены
во~времени (раздел~\ref{sec:tx-dual-single}). Dual-message у~UnionPay появился
позже~--- по~мере роста кредитной эмиссии, международного расширения
и~сценариев с~отложенным подтверждением (гостиницы,
аренда)~\cite{unionpay-russia-rules}.

UnionPay сегодня поддерживает обе схемы. \highlight{В~сценарии
  \textit{offline-single-message} карта принимает транзакцию без
  онлайн-подтверждения и~накапливает офлайн-записи для~последующей
передачи}~\cite{emvco-specs}; такой сценарий применим для~небольших сумм и~ряда
автоматических POS-терминалов. В~схеме \textit{single-message} нет второго
клирингового файла, и~коннектор эквайера обрабатывает транзакцию иначе, чем
в~\textit{dual-message}.

Сценарии с~отложенным подтверждением продавец сверяет по~продукту и~рынку
отдельно. \textit{Pre-auth} (предварительное резервирование суммы на~карте
  без~её немедленного списания; финальное предъявление выполняет отдельная
операция \textit{completion}) в~гостиничном канале работает как партнёрская
функция. IHG (InterContinental Hotels Group, международная гостиничная
группа) описывает запуск приёма бронирований UnionPay с~предавторизацией
без~предъявления карты для~своих прямых каналов в~Большом
Китае~\cite{ihg-unionpay-preauth-2024}; параметры \textit{pre-auth},
\textit{completion} и~последующих изменений суммы у~каждого партнёра свои.

\section{QuickPass: ядро и~протокольный стек}\label{sec:up-quickpass}

Официальные материалы UPI определяют QuickPass как сервис на~базе NFC
(near-field communication~--- бесконтактная связь малой дальности
на~13{,}56~МГц) и~токенизации, который поддерживает бесконтактную
офлайн-оплату, оплату по~QR и~платежи внутри приложения, включая мобильные
телефоны, HCE (host card emulation~--- программная эмуляция карты в~ОС
телефона) и~носимые устройства, если они соответствуют спецификациям
EMVCo (отраслевой консорциум~--- владелец спецификаций EMV: контактный/бесконтактный
чип, токенизация, 3-D Secure~\cite{emvco-specs})~\cite{unionpay-quickpass-overview}.

Физический уровень здесь тот~же, что у~Visa payWave и~Mastercard
PayPass\footnote{Visa payWave и~Mastercard PayPass~--- исторические
  маркетинговые бренды; с~2016~года Mastercard перевела продукт на~бренд
  \enquote{Mastercard Contactless} (ядро M/Chip Contactless), а~Visa к~концу
  2010-х в~маркетинге перешла на~термин \enquote{Visa Contactless}. В~технической
  документации и~сертификационных материалах исторические названия продолжают
использоваться.}: ISO/IEC~14443 (стандарт \textit{proximity}-карт и~физический
уровень NFC на~13{,}56~МГц)~\cite{iso-14443-1}. Один считыватель на~аппаратном
уровне обслуживает несколько бесконтактных сетей, если в~нём установлены
соответствующие ядра приложений (kernels~--- программные компоненты терминала,
  реализующие правила EMV-обработки для~каждой платёжной сети
бесконтактной транзакции).

У~каждой платёжной сети собственное ядро бесконтактной обработки. У~Visa
это VCPS (Visa Contactless Payment Specification), у~Mastercard~---
PayPass/M/Chip Contactless, у~UnionPay~--- отдельное QuickPass-ядро. Каждое
ядро реализует правила своей сети для~выбора приложения (AID,
  application identifier~--- идентификатор EMV-приложения на~карте;
см.~гл.~\ref{ch:emv}), набор команд APDU (application protocol data unit~---
  формат команд и~ответов между терминалом и~чипом
по~ISO/IEC~7816-4~\cite{iso-7816-4}), криптографические параметры и~поведение
терминала при~различных условиях транзакции. Терминал с~одним только ядром VCPS
обнаружит карту UnionPay на~физическом уровне, но не~проведёт транзакцию:
логики обмена с~этим приложением у~него нет.

Выбор приложения EMV через список AID на~карте и~в~терминале (механика
и~определение AID~--- раздел~\ref{sec:emv-app-selection}) определяет,
может~ли транзакция начаться. Карта UnionPay содержит собственный AID
из~диапазона, зарегистрированного за~China UnionPay. Если в~терминальной
конфигурации нет соответствующего AID с~привязанным QuickPass-ядром, терминал
не выберет приложение и~транзакция не начнётся.

UPI публикует отдельные списки рекомендованных QuickPass-точек и~указывает, что
фактический приём отличается от~списков и~сверяется по~знаку
QuickPass в~точке оплаты~\cite{unionpay-quickpass-merchants}. Эквайер
и~интегратор проверяют поддержку бесконтактного интерфейса и~наличие
QuickPass-ядра в~терминальной конфигурации.

Сценарий QR в~QuickPass устроен иначе, чем NFC: обе схемы~--- MPM
(merchant-presented mode~--- QR показывает продавец, сканирует покупатель)
и~CPM (consumer-presented mode~--- QR показывает покупатель, сканирует
продавец)~--- работают на~прикладном уровне, не через EMV APDU (таксономия QR
и~различия между MPM и~CPM~--- раздел~\ref{sec:qr-models}). Подключение
QuickPass~QR~--- отдельная интеграция через API, не связанная с~настройкой ядра
NFC на~терминале.

\section{Карточные продукты и~interchange}\label{sec:up-products}

\highlight{В~отличие от~Visa и~Mastercard, UnionPay не публикует единой
мировой таблицы interchange}. Для~части рынков UPI публикует таблицы
в~открытом доступе на~странице Interchange Rate Fees: на~момент сверки в~апреле
2026~года это EEA (European Economic Area~--- Европейская экономическая зона),
Великобритания, Швейцария, Япония, Канада, Малайзия, Казахстан и~ряд других
юрисдикций; актуальный перечень рынков сверяется по~странице
IRF~\cite{unionpay-irf}. Для~остальных юрисдикций, в~том числе для~самого
Китая, условия определяются договорами с~участниками и~зависят от~типа
продукта, географии и~роли участника в~расчётной системе. На~рынке
без~публичной таблицы условия запрашиваются через договорной процесс
с~участником.

Продуктовая линейка UnionPay структурирована по~уровням: UnionPay Standard,
UnionPay Gold, UnionPay Platinum, UnionPay Diamond и~несколько
специализированных продуктов (UnionPay Business Card, UnionPay Prepaid). Как
у~Visa и~Mastercard, тип продукта влияет на~маршрутизацию, на~правила
управления рисками и~(там, где условия определены) на~стоимость обработки.

UnionPay появился как сеть с~тесной связью между эмитентом, процессингом
и~расчётом. Внутри Китая большинство транзакций проходит через внутреннюю
инфраструктуру UnionPay напрямую, без промежуточных звеньев. За~пределами Китая
операции проходят через инфраструктуру UPI, и~до~подключения команда определяет,
кто именно участник и~какая редакция правил действует в~конкретной юрисдикции.

ТСП, принимающее карты UnionPay через интернет, аутентифицирует транзакции без
физического присутствия карты (\textit{card-not-present}, CNP) одним из~двух
механизмов UPI: UnionPay~3DS, сертифицированный по~спецификации
EMVCo~3DS~\cite{unionpay-3ds-emvco}, или UPI SecurePlus~--- протокол поверх API
с~SMS-аутентификацией и~токенизацией, доступный за~пределами материкового
Китая~\cite{unionpay-secureplus}.\footnote{SecurePlus пришёл на~смену более
  раннему UPI~SecurePay с~HPP (\textit{hosted payment page}) и~редиректом
  покупателя на~страницу UnionPay; SecurePlus встроен в~чекаут эквайера и~не
требует редиректа~\cite{unionpay-secureplus}.} Если эмитент поддерживает один
из~протоколов, ответственность в~случае спора переносится (логика та~же, что
в~разделе~\ref{sec:3ds-liability}); оба протокола относятся к~UnionPay
и~требуют отдельной интеграции.

\section{UnionPay International и~правила участия}\label{sec:up-international}

За~пределами материкового Китая карту обслуживает UnionPay International.
В~собственном профиле компания называет себя дочерней структурой China
UnionPay, отвечающей за~международный бизнес; участие оформляется через
партнёрские институты~\cite{unionpay-international-about}.

Опубликованные правила указывают, что местные правила UnionPay построены
на~актуальной редакции UnionPay International Operating Regulations с~учётом
местного законодательства~\cite{unionpay-russia-rules}. Те~же правила задают
критерии участия для~прямых и~косвенных участников, обязанности эквайера
и~эмитента, требования к~управлению рисками и~информационному
обмену~\cite{unionpay-russia-rules}.

Подключение к~UnionPay начинается с~описания модели участия: через какого
участника принимаются операции, какие договоры и~лицензии оформляют участие,
кто отвечает за~досье на~ТСП и~управление рисками, как оформлены
расчёты~\cite{unionpay-russia-rules}.

Открытые правила описывают расчётный механизм: нетто-клиринг, передачу реестра
нетто-позиций в~расчётный центр, дебетование и~кредитование счетов участников,
а~также конкретные сроки обработки по~дням~T
и~T+1~\cite{unionpay-russia-rules}.

\section{Расчётная валюта и~конвертация}\label{sec:up-currency}

Внутри Китая расчёты между участниками внутренней инфраструктуры UnionPay
проводятся в~юанях (CNY~--- буквенный код китайского юаня
по~ISO~4217~\cite{iso-4217}). За~пределами Китая UPI рассчитывает участников
в~CNY или в~согласованной валюте расчётной системы; выбор валюты зависит
от~двусторонних договорённостей с~UPI и~региональной редакции правил
участия~\cite{unionpay-russia-rules}.

Карточная транзакция проходит через несколько конверсионных шагов на~разных
этапах. Покупатель платит в~местной валюте (в~рублях или тайских батах),
эквайер выставляет транзакцию в~валюте транзакции, а~расчёт между участниками
выполняется в~расчётной валюте системы с~конвертацией по~курсу момента
клиринга, а~не момента авторизации. Курс зачисления продавцу отличается
от~курса на~терминале.

Общий механизм DCC (dynamic currency conversion~--- конвертация на~терминале
  по~предложенному эквайером или его DCC-провайдером курсу с~раскрытием суммы
покупателю) разобран в~разделе~\ref{sec:mc-dcc}~\cite{mc-dcc-guide}.
По~UnionPay механика та~же, но китайский держатель карты ожидает сумму в~CNY,
и~эквайер показывает её через стороннего провайдера DCC в~обход официального
канала UPI; курс DCC по~картам UnionPay менее выгоден, чем сетевой.
Для~продавца, работающего с~китайскими туристами, это типичный источник жалоб
и~спорных операций.

Отдельный случай~--- кросс-валютные расчёты по~двусторонним соглашениям
UPI с~центральными банками. В~ряде юрисдикций UPI согласовал расчёты
в~валюте страны без промежуточной конвертации через CNY. Условия сверяются
по~актуальной редакции правил участия~\cite{unionpay-russia-rules}.

\section{Кобейджинг Мир/UnionPay: маршрутизация и~выбор банка}\label{sec:up-mir}

UPI и~НСПК в~официальных материалах совместно продвигают выпуск дебетовых карт
Мир--UnionPay банками-участниками НСПК~\cite{unionpay-mir-cobadge}; такая карта
обрабатывается через сеть Мир в~России и~через сеть UnionPay за~её
пределами~\cite{unionpay-mir-cobadge}. Действующие правила UnionPay в~России
дополнительно уточняют: для~эмиссии совместной карты эмитент должен участвовать
в~обеих платёжных системах, а~операции, маршрутизируемые через платёжную систему
UnionPay, регулируются её правилами~\cite{unionpay-russia-rules}.

Кобейджинговая карта содержит два приложения: приложение Мир с~AID национального
диапазона и~приложение UnionPay с~AID китайского диапазона. Терминал при~выборе
приложения получает оба AID и~применяет приоритет из~конфигурации. В~России
терминальная конфигурация отдаёт приоритет Мир; за~рубежом доступно только
приложение UnionPay либо оба, в~зависимости от~настроек эквайера и~конфигурации
терминала.

Маршрут кобейджинговой карты выбирается не один~раз при~выпуске, а~при~каждой
транзакции. Сначала терминал формирует список взаимно поддерживаемых AID
(пересечение списка карты и~списка терминала). Затем эквайер или хост
переопределяет приоритет через параметры TMS (terminal management system~---
  система централизованного управления парком терминалов: их конфигурацией,
ключами и~прошивками), например выставляет приоритет UnionPay выше Мир
на~международных терминалах. Если эмитент настроил маршрутизацию по~типу
операции или по~сумме, разные транзакции одной и~той~же карты маршрутизируются
в~разные сети.

Кобейджинговую карту тестируют по~полной комбинации \enquote{карта $\times$
терминал $\times$ страна $\times$ канал}. Авторизацию за~рубежом выполняет
банк-эмитент через инфраструктуру UnionPay (минуя НСПК); решение о~риске
и~логику отказа формируют правила UnionPay.

Если бизнес эквайера или ТСП рассчитывает на~карту \enquote{Мир--UnionPay} как
на~универсальный международный способ оплаты, доступность маршрута проверяется
по~странам, банкам и~точкам приёма~\cite{unionpay-mir-cobadge}.

\section{Модели отказа при~кобейджинге}\label{sec:up-cobadge-decline}

Кобейджинговая карта добавляет к~стандартным причинам отказа отдельный класс:
отказ из-за разрыва маршрутизации, помимо отказа по~состоянию счёта держателя
карты. Для~команды поддержки и~аналитики жалоба \enquote{карта не~прошла
за~рубежом} разбирается на~четыре сценария.

Первая модель: \textbf{AID не найден}. Терминал не содержит QuickPass-ядра или
AID UnionPay не входит в~конфигурацию, и~приложение не выбрано. Держатель карты
видит отказ ещё до~авторизационного запроса. Хост такой транзакции не видит;
держатель воспринимает это как \enquote{карта не работает}. Диагноз: проблема
на~терминальной стороне эквайера, не на~стороне эмитента.

Вторая модель: \textbf{авторизационный отказ по~коду}. Запрос поступил эмитенту
через сеть UnionPay, но эмитент вернул отказ. Используются те~же коды ISO~8583
(стандарт формата сообщений карточных авторизаций; см.~\cite{iso-8583-2023}),
что и~в~других сетях: 05 (Do not honor), 57 (Transaction not permitted to
cardholder), 62 (Restricted card). Типичная причина кода~62 у~карт UnionPay~---
эмитент не~включил этот диапазон BIN (bank identification number~--- первые
цифры PAN, идентифицирующие эмитента и~продуктовый диапазон) или тип операции
в~разрешённый международный маршрут. Карта физически выпущена, но не~настроена
на~работу через UPI.

Третья модель: \textbf{недоступность маршрута}. Сеть UnionPay между эквайером
и~эмитентом недоступна: нет договора между эквайером и~UPI, нет
корреспондентских связей, или канал временно недоступен. Терминал получает
таймаут или технический отказ (96, 91). Если терминал поддерживает
\textit{fallback} (резервный сценарий обработки~--- переключение
  на~альтернативное приложение, сеть или способ ввода, когда основной путь
недоступен) на~альтернативную сеть и~на~карте есть другое приложение, например
национальная сеть страны эквайера, \textit{fallback} возможен. Для~Мир/UnionPay
за~рубежом \textit{fallback} на~Мир невозможен: инфраструктуры Мир за~пределами
России нет.

Четвёртая модель: \textbf{онлайн-канал (транзакция CNP)}. При интернет-платеже
нет терминала и~выбора AID; процессинг решает, по~какому маршруту направить
транзакцию, опираясь на~BIN и~договорную конфигурацию. Если
маршрут UnionPay не настроен для~этого BIN, транзакция либо отклоняется с~кодом
\enquote{неверный маршрут}, либо не создаётся: шлюз отбрасывает карту на~этапе
валидации формы ввода.

\practicenote{%
  При диагностике отказов кобейджинговых карт проверяют
  по~порядку: был~ли авторизационный запрос (или отказ
  на~стороне терминала), какой код ответа вернул эмитент, направлена~ли
  транзакция через UnionPay или через Мир, и~соответствует~ли
  конфигурация BIN на~стороне эквайера ожидаемому маршруту.
}

\section{Санкции и~операционные риски}\label{sec:up-sanctions}

Санкционные ограничения определяют, в~каких юрисдикциях и~при~каких условиях
карта UnionPay фактически работает, независимо от~того, есть~ли технический
маршрут.

\textbf{Позиция UnionPay.} После февраля 2022~года UnionPay International
приостановил переговоры о~новых продуктах с~российскими банками, а~ряду банков
отказано в~подключении к~сети UPI\footnote{Официальных пресс-релизов UPI
  по~этому вопросу в~открытом доступе опубликовано не было; описание поведения
  системы основано на~отраслевых сообщениях 2022~года и~на~наблюдаемом поведении
сети.}. Отраслевые источники называли причиной риск вторичных санкций США
и~Евросоюза. UPI не американская компания: она напрямую не подпадает под санкции
OFAC (Office of Foreign Assets Control~--- управление Министерства финансов США,
администрирующее санкционные программы) и~не включена в~список SDN (Specially
  Designated Nationals and Blocked Persons List~--- санкционный список
заблокированных лиц и~организаций), но вторичные санкции создают риск для~любого
бизнеса, поддерживающего транзакции с~лицами или организациями из~списка
SDN~\cite{ofac-sanctions-list-service}.

\textbf{Вторичные санкции как механизм.} Вторичные санкции США запрещают
компаниям из~третьих стран вести существенные операции с~определёнными
секторами российской экономики, в~частности с~банками, включёнными в~список SDN
или в~Sectoral Sanctions Identifications list (SSI). Для~UnionPay это означает,
что даже при~технически возможной транзакции её проведение через инфраструктуру
UPI создаёт риск вторичных санкций, включая отключение от~долларовых
корреспондентских счетов и~американских рынков капитала. Прямой запрет
отсутствует~\cite{ofac-compliance-framework}; риск создаёт именно вторичный
механизм.

\textbf{Наблюдаемая картина.} Часть российских банков сохранила выпуск
и~обслуживание карт UnionPay; это банки, не попавшие под~SDN и~не входящие
в~секторальные санкционные категории. По~сообщениям деловых СМИ карты работают
в~странах, не~присоединившихся к~западным санкционным пакетам (Китай, ОАЭ,
Турция, страны АСЕАН, часть африканских рынков); UPI не публиковала официальный
перечень поддерживаемых стран. В~государствах, применяющих санкции ЕС
и~США, эквайер либо блокирует диапазоны BIN российских карт на~входе, либо
получает отказ от~банков-корреспондентов при~попытке рассчитать операцию.

\textbf{Операционный риск для~эмитента.} Для~российского банка, выпускающего
карты UnionPay, основной риск договорной: UPI может в~одностороннем порядке
приостановить или прекратить участие банка в~сети, если санкционный периметр
вокруг банка изменится. Обещание клиенту \enquote{карта работает за~рубежом}
создаёт операционный риск, который техническими мерами не устраняется; снизить
его помогают мониторинг санкционного статуса и~прямое предупреждение держателя
карты об~ограничениях.

\textbf{Операционный риск для~ТСП и~эквайера.} ТСП или PSP (payment service
  provider~--- провайдер платёжных услуг, агрегирующий приём для~ТСП на~стороне
эквайера), подключающий UnionPay, обязан проводить скрининг транзакций
по~санкционным спискам OFAC на~общих основаниях. Поскольку значительная доля
держателей карт~--- граждане КНР, политическое напряжение между США и~Китаем
сохраняется долгосрочным фактором риска.

\paragraph{Проверка по~первоисточнику}
Санкционный периметр вокруг UnionPay меняется быстрее, чем печатная книга
успевает его фиксировать. Актуальный список SDN и~SSI поддерживается
OFAC~\cite{ofac-sanctions-list-service}; позиция конкретного банка проверяется
по~текущей редакции санкционных листов; состояние на~дату публикации этой книги
ориентиром не служит.

\section{Отличия от~Visa и~Mastercard}\label{sec:up-vs-global}

При общей карточной модели UnionPay расходится с~Visa и~Mastercard
по~процессинговой инфраструктуре, формату клиринга, бесконтактному ядру,
токенизации, интернет-аутентификации, выплатам, обновлению реквизитов
и~публикации interchange. Таблица~\ref{tab:up-vs-visa-mc} сводит
трёхстороннее сопоставление; на~неё ссылается и~глава о~Mastercard
(раздел~\ref{sec:mc-vs-visa}) за~разбором пары Mastercard и~Visa.

\begin{table}[htb]
  \begingroup
  \footnotesize
  \setlength{\tabcolsep}{4pt}
  \renewcommand{\arraystretch}{1.2}
  \begin{tabularx}{\textwidth}{
      @{}P{0.22\textwidth}
      Y
      Y
      Y@{}
    }
    \toprule
    \headrow
    \thd{Компонент} & \thd{UnionPay}
    & \thd{Visa} & \thd{Mastercard} \\
    \midrule
    Процессинговая инфр. & внутренняя инфраструктура UnionPay (Китай),
    UPI (международно) & VisaNet (BASE~I + BASE~II)
    & Единый периметр (Authorization Platform + GCMS) \\
    Клиринговый формат & Реестры через инфраструктуру UnionPay/UPI
    & TC33 через BASE~II & IPM через GCMS \\
    Бесконтактное ядро & QuickPass-ядро (отдельное для~UnionPay)
    & VCPS & PayPass/M/Chip Contactless \\
    Токенизация & UPI Token Service
    & VTS + Token Assurance Level & MDES + DSRP \\
    Интернет-аутентификация & UnionPay~3DS (EMVCo) + SecurePlus (SMS/токен)
    & Visa Secure (3DS) & Mastercard Identity Check (3DS) \\
    Выплаты & OCT через UPI (отдельного бренда нет)
    & Visa Direct & Mastercard Send \\
    Обновление реквизитов & Нет публичного аналога VAU/ABU
    & Visa Account Updater (VAU) & Automatic Billing Updater (ABU) \\
    Таблицы interchange & Публикуются для~части рынков (EEA, UK,
    JP, CA и~др.~--- двухбуквенные коды стран по~ISO~3166-1~\cite{iso-3166-1});
    для~остальных договорные
    & Публикуются по~регионам & Публикуются по~регионам \\
    \bottomrule
  \end{tabularx}%
  \endgroup
  \caption{Архитектурные отличия UnionPay, Visa и~Mastercard.}\label{tab:up-vs-visa-mc}
\end{table}

Бесконтактное ядро: терминал с~VCPS или PayPass требует отдельного
QuickPass-ядра для~приёма карт UnionPay, даже если физический уровень NFC
одинаковый.

Закрытые условия interchange: команда, которая строила ценообразование
по~публичным таблицам Visa и~Mastercard, не сможет применить ту~же методику
к~UnionPay~--- условия она получает через договорной процесс
с~участником~\cite{unionpay-russia-rules}.

\practicenote{%
  При добавлении UnionPay в~таблицу маршрутизации заранее разделите
  три вопроса: редакция правил и~участник/расчётная система в~вашей
  юрисдикции; поддерживаемый продукт и~маршрут кобейджа (QuickPass,
  QR, Мир--UnionPay); отдельно~--- интернет-аутентификация через UPI
  SecurePlus.
}
